\section{Introduzione}

Il disturbo da deficit di attenzione e iperattività
(\textit{Attention-Deficit Hyperactivity Disorder}, ADHD) è un
disturbo del neurosviluppo caratterizzato da livelli persistenti di
disattenzione, impulsività e iperattività. Poiché i sintomi
riguardano anche la capacità di selezionare gli stimoli rilevanti,
mantenere l'attenzione e controllare la risposta comportamentale, lo
studio dell'attività cerebrale durante compiti attentivi può fornire
informazioni utili sui meccanismi neurofisiologici associati al
disturbo.

L'elettroencefalografia (EEG) è una tecnica non invasiva che misura
le differenze di potenziale elettrico registrate da elettrodi
posizionati sullo scalpo. Il segnale rilevato da ciascun elettrodo non
corrisponde all'attività di un singolo neurone o di una singola area
cerebrale, ma riflette la sovrapposizione dell'attività sincronizzata
di popolazioni neuronali e la sua propagazione attraverso i tessuti
della testa. Il principale vantaggio dell'EEG è l'elevata risoluzione
temporale, che permette di descrivere variazioni dell'attività
cerebrale nell'ordine dei millisecondi e lo rende particolarmente
adatto allo studio di processi dinamici come l'attenzione.

Le registrazioni EEG multicanale possono essere analizzate attraverso
rappresentazioni differenti e complementari. Una prima possibilità
consiste nello studiare il contenuto spettrale del segnale, cioè come
la sua potenza sia distribuita nelle diverse frequenze. L'attività EEG
viene comunemente suddivisa in bande, tra cui delta, theta, alpha e
beta. Queste bande non corrispondono in modo univoco a una specifica
funzione cognitiva: il loro significato dipende dalla regione
considerata, dallo stato del soggetto e dalla condizione sperimentale.
Per esempio, un aumento dell'attività theta può essere associato a
sonnolenza o riduzione della vigilanza, ma la theta frontale può anche
riflettere processi di controllo cognitivo e monitoraggio. Analogamente,
l'attività alpha può diminuire nelle aree coinvolte nell'elaborazione
visiva oppure aumentare nelle regioni che contribuiscono all'inibizione
di informazioni irrilevanti.

Un indice particolarmente studiato nella letteratura sull'ADHD è il
\textit{theta/beta ratio} (TBR), definito come il rapporto tra la
potenza della banda theta e quella della banda beta:

\begin{equation}
	\mathrm{TBR} =
	\frac{P_{\theta}}{P_{\beta}}.
\end{equation}

Storicamente, diversi studi condotti su bambini con ADHD avevano
riportato un aumento dell'attività theta e, in alcuni casi, una
riduzione dell'attività beta. Un TBR elevato era stato quindi
interpretato come un maggiore peso dell'attività lenta rispetto a
quella più rapida e proposto come possibile indicatore di una
regolazione attentiva meno efficiente. Alcuni dei primi lavori
riportavano una separazione molto marcata tra soggetti con ADHD e
controlli, alimentando l'ipotesi che il TBR potesse essere utilizzato
come biomarcatore diagnostico quantitativo.

La validità di questa ipotesi è stata riesaminata da Arns, Conners e
Kraemer attraverso una meta-analisi dedicata al TBR
\cite{Arns2013}. Per rendere gli studi maggiormente confrontabili, gli autori hanno
considerato esclusivamente il theta/beta ratio calcolato sul segnale
registrato dal canale Cz durante la condizione a occhi aperti, in
bambini e adolescenti di età compresa tra 6 e 18 anni. Cz è
l'elettrodo situato sulla linea mediana nella regione centrale dello
scalpo secondo il sistema internazionale 10--20. La scelta di un'unica
posizione e di una stessa condizione di registrazione ha permesso di
ridurre la variabilità metodologica tra gli studi. Sono stati
inclusi nove studi, per un totale di 1253 soggetti con ADHD e 517
controlli. L'effect size medio, calcolato come Hedges' \(d\), risultava
pari a 0.75 nella fascia 6--13 anni e a 0.62 nella fascia 6--18 anni.
Tali valori sembrerebbero indicare una differenza moderata o elevata;
tuttavia, in entrambe le analisi il test di eterogeneità risultava
significativo. Le differenze tra gli studi erano quindi troppo ampie
per essere attribuite esclusivamente alla variabilità casuale del
campionamento, indicando che il TBR non rappresentava un effetto
stabile e uniforme.

L'analisi degli autori ha inoltre evidenziato una progressiva
riduzione dell'effect size negli studi più recenti. Questa tendenza
non era principalmente determinata da una diminuzione del TBR nel
gruppo ADHD, che rimaneva relativamente stabile, ma da un aumento del
TBR nei gruppi di controllo. I due gruppi diventavano quindi
progressivamente più simili. Tra le possibili spiegazioni vengono
discusse differenze nella selezione dei controlli, nelle condizioni di
registrazione, negli strumenti e nelle procedure di preprocessing,
oltre a fattori fisiologici come sonno, stanchezza e livello di
vigilanza.

La conclusione della meta-analisi non è che il TBR sia completamente
privo di informazione. Un eccesso di theta o un TBR elevato può
caratterizzare  soggetti con ADHD e può avere valore
prognostico, per esempio nella previsione della risposta a specifici
trattamenti. Il suo impiego diagnostico richiederebbe invece una
differenza stabile e riproducibile nella maggior parte dei soggetti,
condizione che non risulta supportata dall'eterogeneità osservata.
Il TBR non può quindi essere considerato, da solo, un marker
diagnostico universale dell'ADHD.

Questa distinzione è particolarmente importante per il presente
lavoro. Il dataset utilizzato è stato acquisito durante un compito di
attenzione visiva, mentre la meta-analisi di Arns et al. considera
registrazioni a occhi aperti in condizioni standardizzate e
limitatamente all'elettrodo Cz. L'attività theta e beta può essere
modulata dal compito e non è quindi possibile assumere a priori che il
TBR debba presentare nel dataset analizzato la stessa direzione
descritta nei primi studi a riposo. Il TBR viene pertanto valutato sia
come insieme di feature autonomo sia all'interno di una
rappresentazione spettrale più ampia, comprendente le potenze delle
diverse bande e la loro distribuzione sui canali e sulle regioni dello
scalpo.

Una rappresentazione complementare del segnale è fornita dalla
\textit{functional connectivity}. Mentre le feature spettrali
descrivono quanta attività sia presente in una determinata banda e in
una determinata posizione, la connettività funzionale descrive il
grado di dipendenza statistica tra segnali registrati in posizioni
diverse. Essa non dimostra l'esistenza di una connessione anatomica
diretta e non consente, di per sé, di stabilire una relazione causale
o una direzione dell'interazione. Permette però di rappresentare
l'EEG come una rete di relazioni tra elettrodi o regioni, anziché come
un insieme di canali indipendenti.

Un riferimento importante per questa seconda linea di analisi è il
lavoro di Ahmadi Moghadam et al. \cite{Moghadam2024}, particolarmente
rilevante perché utilizza lo stesso dataset considerato nel presente
progetto. Gli autori analizzano registrazioni EEG di 121 bambini,
61 con ADHD e 60 controlli, acquisite mediante 19 canali a una
frequenza di campionamento di 128 Hz durante un compito di attenzione
visiva. Dopo il preprocessing, i segnali vengono suddivisi nelle
bande delta, theta, alpha, beta e gamma e segmentati in epoche non
sovrapposte di sei secondi.

Per ogni banda e per ogni epoca, gli autori costruiscono due matrici
di connettività \(19 \times 19\). La prima utilizza il
\textit{Pearson Correlation Coefficient} (PCC), che misura la
covariazione lineare tra le ampiezze di due segnali. La seconda
utilizza il \textit{Phase-Locking Value} (PLV), che misura invece la
stabilità nel tempo della differenza di fase tra due segnali. Le due
misure descrivono quindi aspetti differenti della relazione: due
segnali possono mostrare una debole correlazione delle ampiezze ma
mantenere una relazione di fase relativamente stabile.

Poiché sia la matrice di Pearson sia quella di PLV sono simmetriche,
i due triangoli di ciascuna matrice conterrebbero informazioni
duplicate. Gli autori sfruttano questa proprietà per costruire una
\textit{Fused Connectivity Map}: il triangolo inferiore contiene i
valori di Pearson, mentre quello superiore contiene i valori di PLV.
Si ottiene così una matrice della stessa dimensione di una singola
mappa di connettività, ma contenente due forme complementari di
informazione.

Le mappe fuse vengono successivamente trattate come immagini e
fornite in ingresso a una rete neurale convoluzionale dotata di un
meccanismo di \textit{channel attention}. I filtri convoluzionali
apprendono configurazioni locali nella matrice di connettività, mentre
il modulo di attention assegna un peso maggiore alle feature map
considerate più informative per la classificazione. In questo
contesto, il termine ``channel'' nel modulo di attention non indica
gli elettrodi EEG, ma le feature map interne prodotte dagli strati
convoluzionali. Il modello migliore riportato dagli autori utilizza
le mappe della banda theta e raggiunge un'accuracy del
98.88\%.

Nel presente progetto non viene replicata integralmente questa
architettura. L'obiettivo non è confrontare direttamente un modello
tradizionale con la rete neurale proposta nel paper, poiché
preprocessing, rappresentazione dei dati, segmentazione e procedura
di addestramento sono differenti. Viene invece ripresa la sua
motivazione centrale: verificare se le relazioni tra segnali
contengano informazione discriminante ulteriore rispetto alle sole
potenze spettrali. Per mantenere la pipeline più semplice e
interpretabile, la connettività viene stimata mediante correlazione
di Pearson e sintetizzata sia a livello di coppie di canali sia a
livello di coppie di regioni. 

Le feature ottenute vengono quindi
valutate mediante Random Forest e Logistic Regression.

L'obiettivo generale del lavoro è dunque confrontare differenti
rappresentazioni dello stesso segnale EEG per la classificazione di
soggetti con ADHD e soggetti di controllo. Non è individuare a priori un singolo biomarcatore dell'ADHD, ma stabilire quale rappresentazione del
segnale conservi l'informazione più utile per distinguere i due
gruppi.